\section{Completion of the proof}\label{sec:conclusion}

\begin{proof}[Proof of Theorem~\ref{thm:main}]
The case $n=1$ was settled in Section~\ref{sec:setup}; assume $n>1$.
By Lemma~\ref{lem:betti-bound}, $b_2(X)\geq4$.
If \(b_2(X)\geq5\), Proposition~\ref{def:large-betti} gives quotient
singularities for \(B\). If \(b_2(X)=4\) and the polarized family of
smooth fibers of the original \(f\) is isotrivial,
Proposition~\ref{iso:exclusion} gives the same conclusion.
In the remaining case, Proposition~\ref{secsys:rank-four-variation}
chooses the two systems of Proposition~\ref{secsys:second-system} on a
projective IHS deformation \(X'\), still with \(b_2(X')=4\):
projective Lagrangian fibrations \(f':X'\to B\) and \(g:X'\to C\), with
\(C\) normal and projective, such that \(g\) has nonzero polarized
variation and \(f'\) restricts to a finite surjection on every smooth
fiber of \(g\).
If \(B\) had a nonquotient germ, a maximal chart would be singular.
Propositions~\ref{str:minimum} and~\ref{str:products} would supply the
fixed stratum \((S_0,\cS)\) and its product presentations.
Proposition~\ref{eu:euler-zero}, applied to the original fibration,
would give \(\chi_{\mathrm{top}}(S_0)=0\).
For the two fibrations on \(X'\), Proposition~\ref{var:exclusion}
instead gives \(\chi_{\mathrm{top}}(S_0)>0\), a contradiction.
Thus every germ of \(B\) is a quotient singularity in this case as well.
The three cases exhaust the possibilities.

Using the fixed-base family, Proposition~\ref{secsys:second-system}
supplies in every case a finite cover of the unchanged \(B\) by an
\(n\)-dimensional abelian variety:
restrict its first fibration to a smooth fiber of the second.
This cover, the quotient singularities just proved, and the original
projective fibration \(f:X\to B\) meet the hypotheses of
Proposition~\ref{quot:smooth-base}. Hence \(B\) is smooth.
The original $X$ is projective, $f$ is surjective with connected
positive-dimensional fibers, and $B$ is now smooth and projective.
Hwang's theorem applies with precisely these hypotheses
\cite[Theorem~1.2, in the setting of Theorem~1.1]{Hwang2008}, giving
$B\simeq\PP^n_{\C}$.
\end{proof}
